\documentclass[11pt,a4paper]{letter}
\usepackage[T1]{fontenc}
\usepackage[a4paper,margin=1in]{geometry}
\usepackage{lmodern}
\begin{document}
\begin{letter}{Editor-in-Chief\\International Journal of Engineering Science}
\opening{Dear Editor,}

Please consider my manuscript entitled \emph{Directional Wave Dispersion and Energy-Flow Steering in Two-Dimensional Strain-Gradient Elastic Solids} for publication as an original research article in the \emph{International Journal of Engineering Science}.

The paper examines wave propagation in a higher-order elastic continuum with directional material length scales and independent gradient micro-inertia. It combines a local Mindlin Form-II constitutive model with a conforming Bogner--Fox--Schmit discretization of the two-dimensional Bloch eigenproblem. The work relates to the journal's coverage of continuum mechanics, advanced material models and wave propagation.

The main contribution is the separation of directional dispersion and energy-flow steering from complete band-gap formation. Rotating an ellipsoidal length-scale tensor changes the direction of acoustic energy transport. In the homogeneous model, the largest computed deviation of group velocity from the wave-vector direction is $2.78^\circ$. The homogeneous solid shows directional spectral separations but no complete band gap. A periodic composite is studied separately and exhibits a gap on the tested meshes; its continuum gap width is not claimed to be converged. Analytical limits, mesh refinement and numerical checks are provided for the main calculations.

These results clarify how anisotropic higher-gradient stiffness and micro-inertia influence elastic-wave dispersion and energy transport, while distinguishing them from gaps caused by periodic material contrast.

Thank you for considering the manuscript.

\closing{Sincerely,\\Vipin Gupta\\Department of Mathematics, Gurugram University\\vipin@gurugramuniversity.ac.in}
\end{letter}
\end{document}
